\documentclass[11pt, a4paper]{article}

\usepackage{amsmath, amssymb}
\usepackage{geometry}
\usepackage{hyperref}
\usepackage{parskip}
\usepackage{booktabs}
\usepackage{listings}
\usepackage{xcolor}

\geometry{margin=2.5cm}

\lstset{
    basicstyle=\ttfamily\small,
    backgroundcolor=\color{gray!10},
    frame=single,
    breaklines=true,
    columns=fullflexible,
}

\title{SMEFT Density Recombination via Importance Reweighting\\
       \large Using Three Separately Trained Normalising Flows}
\author{}
\date{\today}

\begin{document}
\maketitle

\section{What Each Flow Has Learned}

Each normalising flow is trained on a different SMEFT weight component.
Training on weight type $k$ teaches the flow the \emph{normalised} version
of that weighted density:
\begin{align}
    p_\mathrm{sm}(x)   &= \frac{w_\mathrm{sm}(x)}{Z_\mathrm{sm}},   &
    Z_\mathrm{sm}   &= \sum_i w_{\mathrm{sm},i}, \\[4pt]
    p_\mathrm{lin}(x)  &= \frac{w_\mathrm{lin}(x)}{Z_\mathrm{lin}},  &
    Z_\mathrm{lin}  &= \sum_i w_{\mathrm{lin},i}, \\[4pt]
    p_\mathrm{quad}(x) &= \frac{w_\mathrm{quad}(x)}{Z_\mathrm{quad}},&
    Z_\mathrm{quad} &= \sum_i w_{\mathrm{quad},i}.
\end{align}
The three weight components at Wilson coefficient $c_{Ht} = \pm 5$ are defined as:
\begin{align}
    w_\mathrm{sm}   &= \texttt{eventWeight}, \\
    w_\mathrm{lin}  &= \frac{w(+5) - w(-5)}{10}, \\
    w_\mathrm{quad} &= \frac{w(+5) + w(-5) - 2\,w_\mathrm{sm}}{50}.
\end{align}

Each flow supports two operations:
\begin{itemize}
    \item \textbf{Sampling}: decode $z \sim \mathcal{N}(0, I_d)$ to obtain $x \sim p_k(x)$.
    \item \textbf{Density evaluation}: encode $x \mapsto z$ and compute the
          log-likelihood via the change-of-variables formula.
\end{itemize}

\section{The Physical SMEFT Density}

The true \emph{unnormalised} SMEFT event density at Wilson coefficient $c$ is the
weighted sum of the three components:
\begin{equation}
    \rho(x;\, c) = Z_\mathrm{sm}\, p_\mathrm{sm}(x)
                 + c\, Z_\mathrm{lin}\, p_\mathrm{lin}(x)
                 + c^2\, Z_\mathrm{quad}\, p_\mathrm{quad}(x).
\end{equation}
The $Z_k$ factors restore the physical cross-section scaling that was divided out
during training. The correctly normalised version is:
\begin{equation}
    p(x;\, c) = \frac{\rho(x;\, c)}{Z(c)},
    \qquad
    Z(c) = Z_\mathrm{sm} + c\, Z_\mathrm{lin} + c^2\, Z_\mathrm{quad}.
\end{equation}
At $c = 0$ one recovers the pure SM density; at nonzero $c$ all three flows
contribute.

\section{Why Decoding Three Times from the Same Latent Point Fails}

Training flow $k$ finds a bijection $f_k$ such that $f_k(x)$ pushes $p_k(x)$ forward
to $\mathcal{N}(0, I_d)$.
Two different densities $p_\mathrm{sm}$ and $p_\mathrm{lin}$ require two completely
different bijections to achieve this.
The same latent vector $z$ decodes to unrelated points in data space under
$f_\mathrm{sm}^{-1}$ versus $f_\mathrm{lin}^{-1}$ — there is no shared
coordinate system between the three latent spaces.

\section{Density Evaluation in Normalising Flows}

A normalising flow learns a bijection $f_\theta: \mathcal{X} \to \mathcal{Z}$ such that
\begin{equation}
    z = f_\theta(x) \sim \mathcal{N}(0, I_d).
\end{equation}
The inverse transformation $x = f_\theta^{-1}(z)$ is used for sampling.

To \textbf{evaluate} the density $p_\theta(x)$ at a data point $x$, apply the
change-of-variables formula from probability theory:
\begin{equation}
    p_\theta(x) = p_\mathcal{Z}(f_\theta(x)) \cdot \left|\det \frac{\partial f_\theta(x)}{\partial x}\right|.
    \label{eq:cov}
\end{equation}
In log space (numerically more stable):
\begin{equation}
    \log p_\theta(x) = \log p_\mathcal{Z}(f_\theta(x)) + \log\left|\det J_f(x)\right|,
    \label{eq:cov_log}
\end{equation}
where $J_f(x) = \partial f_\theta(x) / \partial x$ is the Jacobian matrix of the encoder.

\subsection{Understanding Each Term}

\paragraph{Term 1: The Latent Density $\log p_\mathcal{Z}(f_\theta(x))$.}

The encoder $f_\theta$ transforms data $x$ into latent space $z = f_\theta(x)$.
Since by design we enforce $f_\theta(x) \sim \mathcal{N}(0, I_d)$, the latent
log-density is simply:
\begin{equation}
    \log p_\mathcal{Z}(z) = -\frac{1}{2}\|z\|^2 - \frac{d}{2}\log(2\pi).
\end{equation}
This is just a standard multivariate Gaussian and costs $O(d)$ to evaluate.

\paragraph{Term 2: The Log-Determinant Jacobian $\log|\det J_f(x)|$.}

The Jacobian determinant accounts for how the transformation $f_\theta$ stretches or compresses volume in data space.
From the change-of-variables theorem in multivariable calculus, if a bijection
compresses a region by a factor of 2, the density in that region must increase
by a factor of 2 to conserve total probability mass.
Formally:
\begin{equation}
    \int_A p_\theta(x) \, dx = \int_{f_\theta(A)} p_\mathcal{Z}(z) \, dz,
\end{equation}
and this equality is maintained only if $p_\theta$ includes the Jacobian determinant.

The normalising flow is specifically designed so this log-determinant is
computationally tractable. Common strategies include:
\begin{itemize}
    \item \textbf{Autoregressive coupling}: The Jacobian is triangular, so
          $\log|\det J| = \sum_{i=1}^d \log|J_{ii}|$ (diagonal sum).
    \item \textbf{Coupling layers}: The Jacobian is block-diagonal, making
          determinant computation a product of smaller determinants.
    \item \textbf{Continuous flows (Neural ODEs)}: The log-determinant is
          computed via the trace of the Jacobian along the ODE trajectory.
\end{itemize}

\subsection{Concrete Example: Evaluating All Three Flows}

For each sample $x_i$ drawn from the SM flow, you evaluate its likelihood under
all three flows:

\paragraph{SM flow encoding:}
\begin{lstlisting}
z_sm = f_sm.encode(x_i)                    # (d,)
log_det_sm = f_sm.log_det_jacobian(x_i)    # scalar
log_p_sm = gaussian_logpdf(z_sm) + log_det_sm
\end{lstlisting}

\paragraph{Linear flow encoding:}
\begin{lstlisting}
z_lin = f_lin.encode(x_i)
log_det_lin = f_lin.log_det_jacobian(x_i)
log_p_lin = gaussian_logpdf(z_lin) + log_det_lin
\end{lstlisting}

\paragraph{Quadratic flow encoding:}
\begin{lstlisting}
z_quad = f_quad.encode(x_i)
log_det_quad = f_quad.log_det_jacobian(x_i)
log_p_quad = gaussian_logpdf(z_quad) + log_det_quad
\end{lstlisting}

You now have three log-densities for the same point $x_i$:
$\log p_\mathrm{sm}(x_i)$, $\log p_\mathrm{lin}(x_i)$, $\log p_\mathrm{quad}(x_i)$.
These are used to compute the density ratios in Step~3 below.

\subsection{Why the Jacobian is Critical}

Without the Jacobian term, you would compute:
\begin{equation}
    \text{(incorrect)} = \log p_\mathcal{Z}(f_\theta(x)).
\end{equation}
This is the \emph{latent space} density, not the data space density.

Consider a toy example: suppose the flow stretches the latent Gaussian by a
factor of 10 in one direction when mapping to data space.
The same amount of probability mass now spreads over a 10× larger volume, so
data space density should be 10× lower per unit volume.
The Jacobian determinant $|\det J| = 1/10$ captures exactly this: the log-density
is reduced by $\log(1/10) \approx -2.3$, properly accounting for the stretching.

Conversely, if the flow compresses data into a small region ($|\det J| > 1$),
the log-density increases, reflecting higher probability concentration.

\subsection{Mathematical Exactness}

The change-of-variables formula Equation~\eqref{eq:cov} is \emph{exact}, not an approximation.
If $f_\theta$ is a true bijection (which by design it is), the densities integrate correctly:
\begin{equation}
    \int_{\mathcal{X}} p_\theta(x) \, dx = \int_{\mathcal{Z}} p_\mathcal{Z}(z) \, dz = 1.
\end{equation}
The computed density $p_\theta(x)$ is the true likelihood of $x$ under the learned distribution.
This is what makes importance reweighting (Section~\ref{sec:importance}) work exactly.

\section{Correct Approach: Importance Reweighting}
\label{sec:importance}

\subsection{Procedure}

\paragraph{Step 1 — Generate a proposal sample.}
Draw $N$ samples from the SM flow, as it carries the largest $Z_k$ and gives
the best coverage of the relevant phase space:
\begin{equation}
    \{x_1, \ldots, x_N\} \sim p_\mathrm{sm}(x).
\end{equation}

\paragraph{Step 2 — Evaluate log-likelihoods under all three flows.}
Pass each sample $x_i$ through the \emph{encoders} of all three flows to obtain:
\begin{equation}
    \log p_\mathrm{sm}(x_i), \quad
    \log p_\mathrm{lin}(x_i), \quad
    \log p_\mathrm{quad}(x_i).
\end{equation}

\paragraph{Step 3 — Compute log density ratios.}
\begin{align}
    \log r^\mathrm{lin}_i  &= \log p_\mathrm{lin}(x_i)  - \log p_\mathrm{sm}(x_i), \\
    \log r^\mathrm{quad}_i &= \log p_\mathrm{quad}(x_i) - \log p_\mathrm{sm}(x_i).
\end{align}
These ratios are computed once and stored; steps 4--5 below can then be evaluated
instantly for any $c$.

\paragraph{Step 4 — Compute importance weights for any $c$.}
The importance weight corrects for the mismatch between the proposal $p_\mathrm{sm}$
and the target density $p(x;\, c)$:
\begin{equation}
    \tilde{w}_i(c)
    = \frac{p(x_i;\, c)}{p_\mathrm{sm}(x_i)}
    = \frac{Z_\mathrm{sm}
            + c\, Z_\mathrm{lin}\, e^{\log r^\mathrm{lin}_i}
            + c^2\, Z_\mathrm{quad}\, e^{\log r^\mathrm{quad}_i}}
           {Z(c)}.
    \label{eq:importance_weight}
\end{equation}

\paragraph{Step 5 — Form the reweighted sample.}
The collection $\{x_i,\, \tilde{w}_i(c)\}$ is a valid (importance-weighted)
representation of $p(x;\, c)$.
Normalised weights are obtained as $\hat{w}_i = \tilde{w}_i / \sum_j \tilde{w}_j$.

\subsection{Pseudocode Summary}

\begin{lstlisting}
# Steps 1-3: done once
x_samples  = sm_flow.sample(N)                  # decode from N(0,I)
log_psm    = sm_flow.log_prob(x_samples)        # encode direction
log_plin   = lin_flow.log_prob(x_samples)
log_pquad  = quad_flow.log_prob(x_samples)

log_r_lin  = log_plin  - log_psm
log_r_quad = log_pquad - log_psm

# Step 4-5: instant for any c
def get_weights(c):
    unnorm = (Z_sm
              + c    * Z_lin  * exp(log_r_lin)
              + c**2 * Z_quad * exp(log_r_quad))
    return unnorm / unnorm.sum()
\end{lstlisting}

\section{Practical Considerations}

\subsection{Coverage and Variance}

Importance sampling degrades when the proposal $p_\mathrm{sm}$ has low support in
regions where the target $p(x;\, c)$ has significant density.
For small $c$ (e.g.\ $|c| \lesssim 1$) near the SM point this is not a concern.
For large $c$, a better proposal would be the full SMEFT density sampled at that $c$.

\subsection{Effective Sample Size}

A standard diagnostic for the quality of the importance weights is the effective
sample size:
\begin{equation}
    N_\mathrm{eff} = \frac{\left(\sum_i \tilde{w}_i\right)^2}{\sum_i \tilde{w}_i^2}.
\end{equation}
If $N_\mathrm{eff} \approx N$ the reweighting is cheap and reliable.
A small $N_\mathrm{eff} \ll N$ indicates weight collapse and requires either more
initial samples or a closer proposal distribution.

\subsection{Non-linear Recombination}

Equation~\eqref{eq:importance_weight} implements the non-linear recombination
described conceptually above.
The combination is polynomial in $c$ but depends, through the density ratios
$r_i^\mathrm{lin}$ and $r_i^\mathrm{quad}$, on the full learned geometry of all
three flows.
This provides the complete $c$-dependence without any retraining: once the
density ratios are stored, the reweighted distribution at any $c$ is a simple
arithmetic operation on the existing sample set.

\end{document}
